\documentclass[12pt,a4paper]{article}

% ------------------------------------------------------------
% Encoding, language, typography
% ------------------------------------------------------------
\usepackage[utf8]{inputenc}
\usepackage[T1]{fontenc}
\usepackage[english]{babel}
\usepackage{lmodern}
\usepackage{microtype}

% ------------------------------------------------------------
% Page layout
% ------------------------------------------------------------
\usepackage[a4paper,margin=2.5cm]{geometry}
\usepackage{setspace}
\onehalfspacing

% ------------------------------------------------------------
% Mathematics
% ------------------------------------------------------------
\usepackage{amsmath}
\usepackage{amssymb}
\usepackage{amsfonts}
\usepackage{mathtools}
\usepackage{bm}
\usepackage{bbm}

% ------------------------------------------------------------
% Tables, lists, links
% ------------------------------------------------------------
\usepackage{booktabs}
\usepackage{enumitem}
\usepackage{xcolor}
\usepackage[hidelinks]{hyperref}

% ------------------------------------------------------------
% Citations
% ------------------------------------------------------------
\usepackage[round,authoryear]{natbib}

% ------------------------------------------------------------
% Theorem-style environments
% ------------------------------------------------------------
\usepackage{amsthm}
\newtheorem{assumption}{Assumption}
\newtheorem{definition}{Definition}
\newtheorem{proposition}{Proposition}
\newtheorem{remark}{Remark}

% ------------------------------------------------------------
% Useful math commands
% ------------------------------------------------------------
\newcommand{\E}{\mathbb{E}}
\newcommand{\En}{\mathbb{E}_N}
\newcommand{\R}{\mathbb{R}}
\newcommand{\ind}{\mathbbm{1}}

\title{Wald AIPW and Double Machine Learning}
\author{}
\date{}

\begin{document}

\maketitle

\section{The Wald-AIPW Estimator}

\subsection{Assumptions}
\label{subsec:wald_aipw_assumptions}

The Wald-AIPW estimator requires a set of assumptions that are strictly stronger than the
unconfoundedness conditions underlying the plain AIPW estimator, because the treatment
$D_i$ may be endogenous and identification now relies on a binary instrument $Z_i$.

\begin{assumption}[Conditional instrument independence]
\label{ass:ci_independence}
$Z_i \perp\!\!\!\perp \{Y_i(d,z),\, D_i(z)\} \mid X_i$ for all $d \in \{0,1\}$, $z \in \{0,1\}$.
\end{assumption}

The instrument is as good as randomly assigned once we condition on the observed
covariates $X_i$. This is the IV analogue of the strong ignorability condition and replaces the
unconfoundedness assumption on the treatment; in contrast to the AIPW setting, no
analogous restriction is placed on $D_i$ directly.

\begin{assumption}[Exclusion restriction]
\label{ass:exclusion}
$Y_i(d,z) = Y_i(d)$ for all $d,z \in \{0,1\}$.
\end{assumption}

The instrument affects the outcome $Y_i$ only through its effect on treatment uptake $D_i$,
not through any direct channel. Combined with Assumption~\ref{ass:ci_independence}, this
ensures that comparing instrument groups isolates a causal effect that runs through $D_i$
alone.

\begin{assumption}[First-stage relevance]
\label{ass:relevance}
$\E[D_i \mid Z_i = 1, X_i] \neq \E[D_i \mid Z_i = 0, X_i]$ almost surely.
\end{assumption}

The instrument must shift treatment probability after conditioning on $X_i$. In the notation
used for estimation below, this requires the first-stage pseudo-outcome $\widetilde{D}^{\,\mathrm{FS}}$
to have a non-zero sample mean, which serves as the denominator of the Wald ratio.

\begin{assumption}[Monotonicity]
\label{ass:monotonicity}
$D_i(1) \geq D_i(0)$ for all $i$.
\end{assumption}

No individual reduces treatment uptake in response to the instrument being switched on.
This rules out defiers. Together with the previous assumptions, monotonicity ensures that
the identified parameter is the local average treatment effect (LATE) for compliers---those
whose treatment status is changed by the instrument---rather than an uninterpretable mixture
\citep[Theorem~1]{angrist1995identification}.

\begin{assumption}[Overlap]
\label{ass:overlap}
$0 < \Pr(Z_i = 1 \mid X_i = x) < 1$ for all $x$ in the support of $X_i$.
\end{assumption}

Both values of the instrument must be possible at every covariate value. This is the IV
analogue of the common-support condition in the unconfoundedness literature and is
necessary to avoid dividing by zero in the propensity-score weights constructed below.

\begin{remark}
Assumptions~\ref{ass:ci_independence}--\ref{ass:overlap} together imply that the LATE
$\tau_{\mathrm{LATE}} = \E[Y_i(1) - Y_i(0) \mid D_i(1) > D_i(0)]$ is nonparametrically
identified from the observed data distribution, without imposing any functional form on
the outcome or treatment equations
\citep[Section~2]{tan2006regression}.
\end{remark}

% ------------------------------------------------------------------
\subsection{From the Wald Ratio to Covariate Adjustment}
\label{subsec:wald_motivation}
% ------------------------------------------------------------------

Under Assumptions~\ref{ass:ci_independence}--\ref{ass:overlap}, the LATE is
nonparametrically identified. Before developing the doubly-robust estimator, it is useful to
trace the logical steps from the simplest possible IV estimator to the fully augmented version,
since each step addresses a concrete limitation of its predecessor.

\paragraph{The canonical Wald estimator.}

In the absence of covariates, or when $Z_i$ is unconditionally independent of potential
outcomes, identification is immediate. The intent-to-treat (ITT) effect of the instrument on
the outcome and the first-stage effect of the instrument on treatment uptake are both
identified as simple differences in group means:
\begin{align}
    \widehat{\mathrm{ITT}}_Y
    &= \En[Y_i \mid Z_i = 1] - \En[Y_i \mid Z_i = 0], \\[4pt]
    \widehat{\mathrm{ITT}}_D
    &= \En[D_i \mid Z_i = 1] - \En[D_i \mid Z_i = 0].
\end{align}
The Wald estimator \citep{wald1940fitting} divides the two:
\begin{align}
    \hat{\tau}_{\mathrm{Wald}}
    =
    \frac{\En[Y_i \mid Z_i = 1] - \En[Y_i \mid Z_i = 0]}
         {\En[D_i \mid Z_i = 1] - \En[D_i \mid Z_i = 0]}.
    \label{eq:wald_estimator}
\end{align}
This estimator is entirely nonparametric in the sense that it imposes no functional form on
the outcome or treatment equations. Its limitation is precisely that it performs no
covariate adjustment. When instrument independence holds only conditionally on $X_i$
(Assumption~\ref{ass:ci_independence}), the raw group means confound the causal effect of
$Z_i$ with differences in covariate distributions across instrument groups, and the Wald
estimator is inconsistent for the LATE.

\paragraph{Why covariate adjustment is necessary.}

The conditional independence assumption states that $Z_i$ is as good as randomly assigned
\emph{within} cells defined by $X_i$, but not marginally. A concrete example is
\citet{card1993using}, where the instrument---proximity to a four-year college---is as good as randomly assigned with respect to potential earnings after controlling for family background and
location characteristics, but unconditionally it is correlated with these same variables
because colleges tend to be located in areas with higher average education levels. Ignoring
$X_i$ would attribute part of the geographical earnings premium to the college-proximity
instrument, biasing the LATE estimate upward.

More formally, because $Z_i \not\perp\!\!\!\perp \{Y_i(d,z)\}$ marginally, the sample split
$\{Z_i = 1\}$ versus $\{Z_i = 0\}$ compares individuals who differ not only in their
instrument value but also in the distribution of $X_i$. Covariate-adjusted variants of the
Wald estimator are therefore needed whenever Assumption~\ref{ass:ci_independence} holds
only conditionally.

\paragraph{Wald-RA: the regression-adjustment benchmark.}

The simplest covariate-adjusted extension replaces the raw group means with outcome
predictions from regression models that condition on both $Z_i$ and $X_i$. Define
\begin{align}
    \hat{Y}(z, X_i) &= \widehat{\E}[Y_i \mid Z_i = z,\, X_i], \\[2pt]
    \hat{D}(z, X_i) &= \widehat{\E}[D_i \mid Z_i = z,\, X_i],
\end{align}
for $z \in \{0,1\}$. The Wald regression-adjustment (Wald-RA) estimator then takes the form
\begin{align}
    \hat{\tau}_{\mathrm{Wald\text{-}RA}}
    =
    \frac{
        \En\bigl[\hat{Y}(1, X_i) - \hat{Y}(0, X_i)\bigr]
    }{
        \En\bigl[\hat{D}(1, X_i) - \hat{D}(0, X_i)\bigr]
    }.
    \label{eq:wald_ra}
\end{align}
Wald-RA is consistent if both outcome models $\hat{Y}(z, \cdot)$ and treatment models
$\hat{D}(z, \cdot)$ are correctly specified. This is a strong requirement: a single
misspecified model---even if the other is correct---introduces first-order bias. In
high-dimensional settings or when the functional form of the outcome or treatment equation
is unknown, this fragility is a serious concern.

\paragraph{The role of propensity score weighting.}

An alternative approach adjusts for confounding not through the outcome model but through
the propensity score of the instrument,
$e(X_i) = \Pr(Z_i = 1 \mid X_i)$.
The key insight introduced by \citet{tan2006regression}, is that reweighting
observations by the inverse of their instrument propensity score can recover the counterfactual
group means without specifying any outcome or treatment model. Specifically, under
Assumptions~\ref{ass:ci_independence} and~\ref{ass:overlap},
\begin{align}
    \E\!\left[\frac{\ind\{Z_i = z\}\, Y_i}{e_z(X_i)}\right]
    =
    \E[Y_i(z, D_i(z))],
    \qquad
    z \in \{0,1\},
    \label{eq:ipw_id}
\end{align}
where $e_1(X_i) = e(X_i)$ and $e_0(X_i) = 1 - e(X_i)$. This identification result follows
by the same argument as in the unconfoundedness AIPW case but with $Z_i$ playing the
role of the treatment \citep[Proposition~1]{tan2006regression}. The resulting
Wald-IPW estimator is consistent if $e(X_i)$ is correctly specified, and is generally
inconsistent if the propensity score model is wrong.

% ------------------------------------------------------------------
\subsection{Double Robustness and the Wald-AIPW Estimator}
\label{subsec:wald_aipw_dr}
% ------------------------------------------------------------------

Both Wald-RA and Wald-IPW achieve consistent estimation under \emph{one} correctly
specified nuisance model: the outcome model for Wald-RA, and the propensity score model
for Wald-IPW. The augmented inverse probability weighting approach of
\citet{tan2006regression} combines the two strategies in a way that inherits consistency
from whichever component is correctly specified---a property known as double robustness.

\paragraph{Identification via AIPW.}

For each instrument value $z \in \{0,1\}$, define the conditional average potential
outcome (CAPO) by
\begin{align}
    \E[Y_i(z,D_i(z)) \mid X_i]
    &= \E[Y_i \mid Z_i = z,\, X_i]
    \;=:\; \mu_Y(z, X_i).
    \label{eq:capo_y}
\end{align}
Similarly, define $\mu_D(z, X_i) = \E[D_i \mid Z_i = z,\, X_i]$.
Under Assumptions~\ref{ass:ci_independence}--\ref{ass:overlap},
the CAPOs can be written in three equivalent ways---as a regression, as an IPW
expression, or as the augmented combination:
\begin{align}
    \E[Y_i(z,D_i(z)) \mid X_i]
    &= \mu_Y(z, X_i)
    \label{eq:id_ra} \\
    &= \E\!\left[
        \frac{\ind\{Z_i = z\}\, Y_i}{e_z(X_i)}
        \;\Big|\; X_i
    \right]
    \label{eq:id_ipw} \\
    &= \E\!\left[
        \mu_Y(z, X_i)
        +
        \frac{\ind\{Z_i = z\}\bigl(Y_i - \mu_Y(z, X_i)\bigr)}{e_z(X_i)}
        \;\Big|\; X_i
    \right].
    \label{eq:id_aipw}
\end{align}
The representation in \eqref{eq:id_aipw} is the doubly-robust form: it equals
$\E[Y_i(z,D_i(z)) \mid X_i]$ whenever either $\mu_Y(z, \cdot)$ or $e(\cdot)$ is correctly
specified, because a misspecification error in one nuisance is multiplied by the residual of
the other \citep[Section~4]{tan2006regression}. The same three representations hold for
$\mu_D(z, X_i)$ with $D_i$ in place of $Y_i$.

\paragraph{What double robustness means.}

To see why the augmented form \eqref{eq:id_aipw} is doubly robust, note that it can be
rewritten as
\begin{align}
    \mu_Y(z, X_i)
    +
    \frac{\ind\{Z_i = z\}\bigl(Y_i - \mu_Y(z, X_i)\bigr)}{e_z(X_i)}
    =
    \mu_Y(z, X_i)
    +
    \underbrace{
        \frac{\ind\{Z_i = z\} - e_z(X_i)}{e_z(X_i)}
    }_{\text{IPW correction weight}}
    \cdot
    \bigl(Y_i - \mu_Y(z, X_i)\bigr).
    \label{eq:dr_decomp}
\end{align}
If the outcome model is correct, $\E[Y_i - \mu_Y(z,X_i) \mid X_i] = 0$, so the IPW
correction term vanishes in expectation regardless of $e(\cdot)$. If instead the propensity
score is correct, $\E[\ind\{Z_i=z\} - e_z(X_i) \mid X_i] = 0$, so the correction term
vanishes regardless of the outcome model. Consistency therefore requires only one of the
two models to be correctly specified, not both simultaneously.

% ------------------------------------------------------------------
\subsection{The Pseudo-Outcome Construction}
\label{subsec:wald_aipw_pseudo}
% ------------------------------------------------------------------

The identification results of Section~\ref{subsec:wald_aipw_dr} directly suggest a
feasible estimator. Let $\hat{\mu}_Y(z, X_i)$, $\hat{\mu}_D(z, X_i)$, and $\hat{e}(X_i)$
denote first-step estimates of the nuisance parameters, and define the instrument-group
inverse probability weights
\begin{align}
    \hat{\lambda}^z_{1,i} = \frac{Z_i}{\hat{e}(X_i)},
    \qquad
    \hat{\lambda}^z_{0,i} = \frac{1 - Z_i}{1 - \hat{e}(X_i)}.
    \label{eq:lambda_z}
\end{align}

\paragraph{Reduced-form pseudo-outcome.}

The AIPW representation in \eqref{eq:id_aipw} applied to the outcome $Y_i$ motivates
the following individual-level pseudo-outcome:
\begin{align}
    \widetilde{Y}^{\,\mathrm{RF}}_i
    &=
    \underbrace{
        \hat{\mu}_Y(1, X_i) - \hat{\mu}_Y(0, X_i)
    }_{\text{outcome model contrast}}
    +
    \underbrace{
        \hat{\lambda}^z_{1,i}
        \bigl(Y_i - \hat{\mu}_Y(1, X_i)\bigr)
        -
        \hat{\lambda}^z_{0,i}
        \bigl(Y_i - \hat{\mu}_Y(0, X_i)\bigr)
    }_{\text{IPW residual correction}}.
    \label{eq:rf_pseudo}
\end{align}
The sample mean $\En[\widetilde{Y}^{\,\mathrm{RF}}_i]$ consistently estimates the
reduced-form effect of $Z_i$ on $Y_i$, i.e.\ the intent-to-treat effect, after covariate
adjustment.

\paragraph{First-stage pseudo-outcome.}

By the exact same construction applied to the treatment variable $D_i$:
\begin{align}
    \widetilde{D}^{\,\mathrm{FS}}_i
    &=
    \underbrace{
        \hat{\mu}_D(1, X_i) - \hat{\mu}_D(0, X_i)
    }_{\text{treatment model contrast}}
    +
    \underbrace{
        \hat{\lambda}^z_{1,i}
        \bigl(D_i - \hat{\mu}_D(1, X_i)\bigr)
        -
        \hat{\lambda}^z_{0,i}
        \bigl(D_i - \hat{\mu}_D(0, X_i)\bigr)
    }_{\text{IPW residual correction}}.
    \label{eq:fs_pseudo}
\end{align}
The sample mean $\En[\widetilde{D}^{\,\mathrm{FS}}_i]$ consistently estimates the
first-stage effect of $Z_i$ on $D_i$.

\paragraph{The Wald-AIPW estimator.}

The Wald-AIPW estimator \citep{tan2006regression} takes the ratio of the two
pseudo-outcome means:
\begin{align}
    \hat{\tau}_{\mathrm{Wald\text{-}AIPW}}
    =
    \frac{
        \En\!\left[\widetilde{Y}^{\,\mathrm{RF}}_i\right]
    }{
        \En\!\left[\widetilde{D}^{\,\mathrm{FS}}_i\right]
    }
    =
    \frac{
        \frac{1}{N}\sum_{i=1}^{N}
        \left[
            \hat{\mu}_Y(1,X_i) - \hat{\mu}_Y(0,X_i)
            + \hat{\lambda}^z_{1,i}(Y_i - \hat{\mu}_Y(1,X_i))
            - \hat{\lambda}^z_{0,i}(Y_i - \hat{\mu}_Y(0,X_i))
        \right]
    }{
        \frac{1}{N}\sum_{i=1}^{N}
        \left[
            \hat{\mu}_D(1,X_i) - \hat{\mu}_D(0,X_i)
            + \hat{\lambda}^z_{1,i}(D_i - \hat{\mu}_D(1,X_i))
            - \hat{\lambda}^z_{0,i}(D_i - \hat{\mu}_D(0,X_i))
        \right]
    }.
    \label{eq:wald_aipw_estimator}
\end{align}

The structure of \eqref{eq:wald_aipw_estimator} makes the parallel to the canonical
Wald ratio in \eqref{eq:wald_estimator} transparent: numerator and denominator both
have the same AIPW form, differing only in whether $Y_i$ or $D_i$ appears as the
response variable. In the absence of covariates, with a known propensity score of
$\hat{e}(X_i) = 1/2$, and with constant outcome and treatment predictions, the
Wald-AIPW reduces exactly to \eqref{eq:wald_estimator}.

\begin{remark}[Double robustness of the Wald-AIPW]
\label{rem:wald_aipw_dr}
The Wald-AIPW estimator in \eqref{eq:wald_aipw_estimator} is doubly robust in the
following sense: both the reduced-form numerator $\En[\widetilde{Y}^{\,\mathrm{RF}}_i]$ and
the first-stage denominator $\En[\widetilde{D}^{\,\mathrm{FS}}_i]$ are individually consistent
if either the respective outcome or treatment model \emph{or} the instrument propensity
score $e(X_i)$ is correctly specified. Consistency of the ratio then requires that both
numerator and denominator converge to their respective population quantities, which holds
if, for each component, at least one of the two nuisance models is correctly specified.
This distinguishes Wald-AIPW from Wald-RA (consistent only if both outcome and
treatment models are correct) and from Wald-IPW (consistent only if the propensity
score is correct).
\end{remark}

% ------------------------------------------------------------------
\subsection{Comparison of the Three Wald Estimator Variants}
\label{subsec:wald_comparison}
% ------------------------------------------------------------------

To fix ideas, Table~\ref{tab:wald_comparison} summarizes the three estimators and the
conditions under which each is consistent.

\begin{table}[h!]
\centering
\caption{Wald estimator variants and their consistency requirements}
\label{tab:wald_comparison}
\smallskip
\begin{tabular}{lccc}
\toprule
\textbf{Estimator} & \textbf{Uses $\hat{\mu}$} & \textbf{Uses $\hat{e}$}
    & \textbf{Consistent if} \\
\midrule
Wald-RA   & \checkmark & $\times$
    & both $\hat{\mu}_Y$ and $\hat{\mu}_D$ correct \\[2pt]
Wald-IPW  & $\times$   & \checkmark
    & $\hat{e}$ correct \\[2pt]
Wald-AIPW & \checkmark & \checkmark
    & $\hat{\mu}$ \emph{or} $\hat{e}$ correct (per component) \\
\bottomrule
\end{tabular}
\end{table}

The double-robustness of Wald-AIPW provides important insurance in practice. When the
functional form of the outcome or treatment model is uncertain---as is typical with a rich
covariate vector $X_i$---the propensity score model for $Z_i$ may be much simpler and more
reliably estimated, since the instrument is often near-randomly assigned by design. Conversely,
in settings with imperfect overlap (propensity scores close to 0 or 1), the outcome model
provides an alternative route to consistent estimation.

Moreover, the AIPW structure confers an additional statistical advantage: the score function
underlying \eqref{eq:wald_aipw_estimator} is semiparametrically efficient, meaning that no
regular estimator can achieve a smaller asymptotic variance when the true data-generating
process belongs to the nonparametric model defined by
Assumptions~\ref{ass:ci_independence}--\ref{ass:overlap}
\citep[Theorem~1]{tan2006distributional}.






% ------------------------------------------------------------------
\section{Inference via Double Machine Learning}
\label{subsec:wald_aipw_dml}
% ------------------------------------------------------------------
 
The estimator in \eqref{eq:wald_aipw_estimator} requires four nuisance functions---
$\mu_Y(1,\cdot)$, $\mu_Y(0,\cdot)$, $\mu_D(1,\cdot)$, $\mu_D(0,\cdot)$---and the
instrument propensity score $e(\cdot)$. \citet{tan2006regression} establishes
consistency under parametric nuisance models, but parametric specifications impose
functional form restrictions on how the instrument shifts outcomes and treatment
probabilities as a function of $X_i$. In applications with many covariates or unknown
interaction structures, such restrictions can produce systematic bias in the nuisance
estimates and, consequently, in the final LATE estimate, even under the double-robustness
property of Section~\ref{subsec:wald_aipw_dr}.
 
The Double Machine Learning (DML) framework of \citet{chernozhukov2018double} resolves
this tension. It shows that the nuisance functions can be estimated with flexible,
nonparametric machine learning methods---random forests, gradient boosting, neural
networks, Lasso---while still obtaining $\sqrt{N}$-consistent, asymptotically normal,
and semiparametrically efficient inference on the LATE, provided two conditions are met:
the estimating score must be \emph{Neyman orthogonal}, and the nuisance estimates must be
obtained by \emph{cross-fitting}. The remainder of this section develops each condition
and explains why both are indispensable.
 
\subsubsection{Why Plug-in Estimation Fails with Machine Learning}
\label{subsubsec:plugin_fails}
 
To understand why standard plug-in estimation breaks down, consider replacing the true
nuisance functions with ML estimates $\hat{\mu}_Y$, $\hat{\mu}_D$, $\hat{e}$ in
\eqref{eq:wald_aipw_estimator} without any further correction. The nuisance estimation
error enters the score as an additive first-order term. A generic ML estimator converges
at a rate slower than $N^{-1/2}$: many methods achieve rates of order $N^{-2/5}$ or
$N^{-1/3}$ under smoothness conditions on the regression function. This means the
contamination from nuisance errors does not vanish fast enough to be absorbed into the
$\sqrt{N}$ asymptotic approximation, and the resulting estimator for the LATE is
asymptotically biased.
 
\citet[p.\ C8]{chernozhukov2018double} make this precise by noting that standard
Donsker-class conditions, which would allow plug-in estimation to work, are violated in
high-dimensional settings where the nuisance parameter spaces grow with $N$---ruling out
even simple linear models with many regressors. The solution is to design the estimating
equation so that nuisance errors enter only in second order.
 
\subsubsection{Neyman Orthogonality}
\label{subsubsec:neyman_orth}
 
\begin{definition}[Neyman orthogonality, \citealp{chernozhukov2018double}]
\label{def:neyman_orth}
A score function $\psi(W_i;\, \tau,\, \eta)$, where $\tau$ is the parameter of interest
and $\eta = (\mu_Y, \mu_D, e)$ collects the nuisance parameters, is \emph{Neyman
orthogonal} at the true values $(\tau_0, \eta_0)$ if the Gateaux derivative of the
expected score with respect to the nuisance parameter is zero:
\begin{align}
    \partial_r\,
    \E\!\left[\psi\!\left(W_i;\, \tau_0,\, \eta_0 + r(\tilde\eta - \eta_0)\right)\right]
    \Big|_{r=0}
    = 0
    \qquad
    \text{for all } \tilde\eta \in \mathcal{T},
    \label{eq:neyman_orth}
\end{align}
where $\mathcal{T}$ is the nuisance realization set.
\end{definition}
 
Condition \eqref{eq:neyman_orth} states that perturbing the nuisance parameters away
from their true values in any direction has no first-order effect on the expected score.
Intuitively, the score is locally insensitive to nuisance misspecification: the bias from
estimating $\eta$ by $\hat\eta$ enters the estimating equation only through
second-order products of nuisance errors, which vanish at the product rate
$\|\hat\eta - \eta_0\|^2$ rather than the level rate $\|\hat\eta - \eta_0\|$.
 
\paragraph{Neyman orthogonality of the Wald-AIPW score.}
 
The Wald-AIPW score can be written compactly as
\begin{align}
    \psi(W_i;\, \tau,\, \eta)
    =
    \underbrace{
        \widetilde{Y}^{\,\mathrm{RF}}_i(\eta)
    }_{\text{reduced form}}
    -
    \tau \cdot
    \underbrace{
        \widetilde{D}^{\,\mathrm{FS}}_i(\eta)
    }_{\text{first stage}},
    \label{eq:wald_aipw_score}
\end{align}
where $\widetilde{Y}^{\,\mathrm{RF}}_i$ and $\widetilde{D}^{\,\mathrm{FS}}_i$ are
defined in \eqref{eq:rf_pseudo} and \eqref{eq:fs_pseudo} respectively.
\citet[Section~5.2, p.~C37]{chernozhukov2018double} verify that this score satisfies
the moment condition $\E[\psi(W_i;\tau_0,\eta_0)]=0$ and the orthogonality condition
\eqref{eq:neyman_orth}.
 
The orthogonality follows from the same argument that establishes double robustness in
identification. A perturbation of the outcome model $\mu_Y(z,\cdot)$ in direction
$\tilde\mu_Y - \mu_Y$ contributes a term proportional to
$\E[(\tilde\mu_Y - \mu_Y)(Z_i - e(X_i))/e(X_i)]$, which is zero by the law of iterated
expectations under Assumption~\ref{ass:ci_independence}:
\begin{align}
    \E\!\left[
        \frac{(\tilde\mu_Y - \mu_Y)(z, X_i)\cdot
        (\ind\{Z_i = z\} - e_z(X_i))}{e_z(X_i)}
    \right]
    =
    \E\!\left[
        (\tilde\mu_Y - \mu_Y)(z, X_i)\cdot
        \E\!\left[
            \frac{\ind\{Z_i = z\} - e_z(X_i)}{e_z(X_i)}
            \,\Big|\, X_i
        \right]
    \right]
    = 0,
    \label{eq:orth_mu}
\end{align}
since $\E[\ind\{Z_i = z\} - e_z(X_i) \mid X_i] = 0$ by definition of the propensity
score. Symmetrically, a perturbation of the propensity score contributes a term
proportional to $\E[(Y_i - \mu_Y(z,X_i))(\tilde e - e)/e^2]$, which is zero because
$\E[Y_i - \mu_Y(z,X_i) \mid X_i] = 0$ at the true outcome model. The same argument
applies to each nuisance in the first-stage component $\widetilde{D}^{\,\mathrm{FS}}_i$.
 
The practical consequence is that if each nuisance function is estimated at a rate $r_N$
(for example $r_N = N^{-2/5}$ for a nonparametric regression estimator under a smoothness
condition), the bias in the final LATE estimate is of order $r_N^2$, not $r_N$. The
condition for valid $\sqrt{N}$ inference is therefore only that
\begin{align}
    r_N^2 = o(N^{-1/2}),
    \qquad \text{i.e.} \qquad
    r_N = o(N^{-1/4}),
    \label{eq:rate_cond}
\end{align}
a much weaker requirement than the $o(N^{-1/2})$ rate that would be needed for direct
plug-in \citep[Assumption~3.2, p.~C25]{chernozhukov2018double}. Random forests and
gradient boosting satisfy \eqref{eq:rate_cond} under mild smoothness or sparsity
conditions on the true nuisance functions, making them suitable for estimating
$\hat\mu_Y$, $\hat\mu_D$, and $\hat e$.
 
\begin{remark}[Comparison with PLR-IV]
\label{rem:orth_vs_plriv}
The Neyman-orthogonal score for the PLR-IV estimator (equation~\eqref{eq:pliv_robinson_score})
also satisfies \eqref{eq:neyman_orth} and requires the same product-rate condition
\eqref{eq:pliv_product_rate}. The Wald-AIPW score in \eqref{eq:wald_aipw_score} imposes
the stronger interactive IV structure of \eqref{ass:ci_independence} rather than the
partially linear outcome equation, allowing for fully nonparametric covariate adjustment
in both the reduced form and the first stage. This is what enables consistent LATE
estimation under unrestricted effect heterogeneity, at the cost of estimating two
additional nuisance functions ($\mu_D(1,\cdot)$ and $\mu_D(0,\cdot)$) relative to
PLR-IV.
\end{remark}
 
\subsubsection{Cross-Fitting}
\label{subsubsec:cross_fitting}
 
Neyman orthogonality alone is not sufficient for valid ML-based inference. Even if
the score is locally insensitive to nuisance errors, the nuisance estimates and the
outcome data must not be from the same sample, or else overfitting creates a
dependence between $\hat\eta$ and the score evaluation that invalidates the central
limit theorem argument.
 
Cross-fitting, introduced formally by \citet[Definitions~3.1--3.2, p.~C23--C24]{chernozhukov2018double},
resolves this by splitting the sample into $K$ folds of roughly equal size $n = N/K$.
For each fold $k \in \{1,\ldots,K\}$, the nuisance functions are estimated on the
complementary $K-1$ folds, producing cross-fitted estimates
$\hat\eta^{(-k)} = (\hat\mu_Y^{(-k)},\, \hat\mu_D^{(-k)},\, \hat e^{(-k)})$ that
are independent of the observations in fold $k$. The pseudo-outcomes are then
constructed using $\hat\eta^{(-k)}$ only for observations in fold $k$:
\begin{align}
    \widetilde{Y}^{\,\mathrm{RF}}_i
    &=
    \hat\mu_Y^{(-k(i))}(1, X_i) - \hat\mu_Y^{(-k(i))}(0, X_i)
    +
    \hat\lambda^{z,(-k(i))}_{1,i}
    \bigl(Y_i - \hat\mu_Y^{(-k(i))}(1,X_i)\bigr)
    -
    \hat\lambda^{z,(-k(i))}_{0,i}
    \bigl(Y_i - \hat\mu_Y^{(-k(i))}(0,X_i)\bigr),
    \label{eq:rf_pseudo_cf}
\end{align}
where $k(i)$ denotes the fold containing observation $i$. The final estimator averages
over all folds:
\begin{align}
    \hat\tau_{\mathrm{Wald\text{-}AIPW}}
    =
    \frac{
        \frac{1}{N}\sum_{k=1}^{K}\sum_{i \in I_k}
        \widetilde{Y}^{\,\mathrm{RF}}_i
    }{
        \frac{1}{N}\sum_{k=1}^{K}\sum_{i \in I_k}
        \widetilde{D}^{\,\mathrm{FS}}_i
    }.
    \label{eq:wald_aipw_cf}
\end{align}
 
Cross-fitting serves two functions. First, the out-of-sample nuisance predictions
$\hat\eta^{(-k)}$ are statistically independent of the fold-$k$ observations, which
allows the empirical process terms in the central limit theorem argument to be controlled
without Donsker-class conditions on the nuisance estimators
\citep[p.~C8]{chernozhukov2018double}. Second, cross-fitting uses the full sample for
both nuisance estimation and score evaluation---unlike a single train/test split, which
would discard half the data---so there is no loss of sample efficiency.
 
\citet[Remark~3.1, p.~C24]{chernozhukov2018double} recommend $K = 4$ or $K = 5$ folds
based on simulation evidence and empirical practice; they note that larger $K$ provides
more training observations for the harder nuisance estimation problem, while the
asymptotic results are valid for any fixed $K \geq 2$.
 
\subsubsection{The Main Inference Result}
\label{subsubsec:wald_aipw_theorem}
 
Combining Neyman orthogonality with cross-fitting yields the central theoretical
justification for using Wald-AIPW with machine learning nuisance estimators.
 
\begin{proposition}[DML inference on the LATE, \citealp{chernozhukov2018double}]
\label{prop:wald_aipw_clt}
Let Assumptions~\ref{ass:ci_independence}--\ref{ass:overlap} hold, and suppose that the
interactive IV model conditions of \citet[Assumption~5.2]{chernozhukov2018double} are
satisfied, including the rate condition
\begin{align}
    \|\hat\mu_Y^{(-k)} - \mu_{Y,0}\|_{P,2}\,
    \|\hat e^{(-k)} - e_0\|_{P,2}
    +
    \|\hat\mu_D^{(-k)} - \mu_{D,0}\|_{P,2}\,
    \|\hat e^{(-k)} - e_0\|_{P,2}
    = o(N^{-1/2}).
    \label{eq:wald_aipw_rate}
\end{align}
Then the DML estimator $\hat\tau_{\mathrm{Wald\text{-}AIPW}}$ in \eqref{eq:wald_aipw_cf}
satisfies
\begin{align}
    \hat\sigma^{-1}\sqrt{N}
    \bigl(\hat\tau_{\mathrm{Wald\text{-}AIPW}} - \tau_{\mathrm{LATE}}\bigr)
    \;\xrightarrow{d}\;
    \mathcal{N}(0,1),
    \label{eq:wald_aipw_clt}
\end{align}
uniformly over $P \in \mathcal{P}$, where
$\hat\sigma^2 = (\En[\widetilde{D}^{\,\mathrm{FS}}_i])^{-2}\,
\En[(\widetilde{Y}^{\,\mathrm{RF}}_i - \hat\tau_{\mathrm{Wald\text{-}AIPW}}\,
\widetilde{D}^{\,\mathrm{FS}}_i)^2]$ is a consistent variance estimator.
Confidence intervals $\hat\tau_{\mathrm{Wald\text{-}AIPW}} \pm
\Phi^{-1}(1-\alpha/2)\hat\sigma/\sqrt{N}$ have uniform asymptotic validity.
\end{proposition}
 
\noindent
The rate condition \eqref{eq:wald_aipw_rate} requires that \emph{products} of nuisance
estimation errors converge at rate $o(N^{-1/2})$. If each of the five nuisance functions
is estimated at a rate $N^{-1/4+\epsilon}$ for some $\epsilon > 0$---a condition satisfied
by random forests under smoothness assumptions \citep[p.~C26]{chernozhukov2018double}---then
the product rate condition is met. This is the operational sense in which Neyman orthogonality
enables the use of machine learning: it demotes the nuisance estimation problem from a
$\sqrt{N}$-rate bottleneck to a nuisance that can be handled by any sufficiently consistent
nonparametric method.
 
The uniformity statement in \eqref{eq:wald_aipw_clt} over $P \in \mathcal{P}$ is
important and distinguishes the DML result from pointwise asymptotic claims.
\citet[p.~C27]{chernozhukov2018double} note that uniformity holds precisely because the
score is Neyman orthogonal: confidence intervals constructed from \eqref{eq:wald_aipw_clt}
remain valid under sequences of data-generating processes that drift within $\mathcal{P}$,
making them robust to local misspecification in a way that is impossible for estimators
based on non-orthogonal scores.
 
\begin{remark}[No bias correction needed]
\label{rem:no_bias_correction}
A direct consequence of Proposition~\ref{prop:wald_aipw_clt} is that no bias-correction
or undersmoothing of the nuisance estimators is required. Standard $t$-tests and
confidence intervals can be applied to $\hat\tau_{\mathrm{Wald\text{-}AIPW}}$ without
modification, because the second-order nuisance bias is $o(N^{-1/2})$ and therefore
negligible relative to the $\sqrt{N}$ fluctuations of the estimator. This contrasts with
direct plug-in estimators, for which the bias is first-order and must be removed
explicitly.
\end{remark}
 
\subsubsection{Summary: Why Wald-AIPW is the Natural DML Estimator for the LATE}
\label{subsubsec:wald_aipw_summary}
 
The properties developed in this section collectively explain why Wald-AIPW, embedded in
the DML framework, is the preferred estimator for the LATE in the presence of rich
covariate information. Table~\ref{tab:wald_aipw_properties} collects the key advantages
relative to the other Wald variants and to the PLR-IV estimator.
 
\begin{table}[h!]
\centering
\caption{Properties of IV estimators used in this thesis}
\label{tab:wald_aipw_properties}
\smallskip
\begin{tabular}{lcccc}
\toprule
\textbf{Property} & \textbf{Wald} & \textbf{Wald-RA} & \textbf{PLR-IV} & \textbf{Wald-AIPW} \\
\midrule
Allows effect heterogeneity  & \checkmark & \checkmark & $\times$ & \checkmark \\
Covariate adjustment         & $\times$   & \checkmark & \checkmark & \checkmark \\
Doubly robust                & $\times$   & $\times$   & $\times$   & \checkmark \\
Neyman orthogonal score      & $\times$   & $\times$   & \checkmark & \checkmark \\
ML nuisance estimation valid & $\times$   & $\times$   & \checkmark & \checkmark \\
Semiparametrically efficient & $\times$   & $\times$   & $\times$   & \checkmark \\
\bottomrule
\end{tabular}
\end{table}
 
The Wald estimator performs no covariate adjustment and is therefore inconsistent under
conditional instrument independence. Wald-RA adjusts for covariates but is fragile to
outcome model misspecification and its score is not Neyman orthogonal, so ML-based
nuisance estimation introduces a first-order bias. PLR-IV has a Neyman-orthogonal score
and supports ML estimation, but requires a partially linear outcome equation and therefore
restricts effect heterogeneity to zero---targeting the structural IV coefficient
$\theta_0$ rather than the LATE. The Wald-AIPW combines covariate adjustment,
double robustness, Neyman orthogonality, and compatibility with arbitrary nonparametric
ML estimators, and under the interactive IV model of \citet{chernozhukov2018double} it
achieves the semiparametric efficiency bound for the LATE. These are the properties that
motivate its use as the DML estimator in this thesis.
 
 
% ------------------------------------------------------------------
\bibliographystyle{apalike}
\bibliography{references}
% ------------------------------------------------------------------
\end{document}